\documentclass[11pt]{article}

% basic packages
\usepackage[margin=1in]{geometry}
\usepackage[pdftex]{graphicx}
\usepackage{amsmath,amssymb,amsthm}
\usepackage{custom}
\usepackage{booktabs}
\usepackage{enumitem}
\usepackage{float}
\usepackage{braket}

% page formatting
\usepackage{fancyhdr}
\pagestyle{fancy}

\renewcommand{\sectionmark}[1]{\markright{\textsf{\arabic{section}. #1}}}
\renewcommand{\subsectionmark}[1]{}
\lhead{\textbf{\thepage} \ \ \nouppercase{\rightmark}}
\chead{}
\rhead{}
\lfoot{}
\cfoot{}
\rfoot{}
\setlength{\headheight}{14pt}

\linespread{1.03} % give a little extra room
\setlength{\parindent}{0.2in} % reduce paragraph indent a bit
\setcounter{secnumdepth}{2} % no numbered subsubsections
\setcounter{tocdepth}{2} % no subsubsections in ToC

\begin{document}

% make title page
\thispagestyle{empty}
\bigskip \
\vspace{0.1cm}

\begin{center}
{\fontsize{22}{22} \selectfont Lecture Notes on}
\vskip 16pt
{\fontsize{36}{36} \selectfont \bf \sffamily Machine Learning Theory}
\vskip 24pt
{\fontsize{18}{18} \selectfont \rmfamily Owen Lin} 
\vskip 6pt
{\fontsize{14}{14} \selectfont \ttfamily owen2006lin@gmail.com} 
\vskip 24pt
\end{center}

{\parindent0pt \baselineskip=15.5pt
These notes are for a graduate level course on machine learning theory, with an 
emphasis on statistical learning theory and online learning.
\begin{itemize}
    \item Shalev-Schwartz and Ben-David: \textit{Understanding Machine Learning: From Theory to Algorithms}
    \item Blum, Hopcroft, Kannan: \textit{Foundations of Data Science}
\end{itemize}

}
 


% make table of contents
\newpage
\microtoc
\newpage

% main content
\section{Dunno}
\subsection{Premise of Learning}
To begin, we'll present some illustrative examples. Suppose we desire to train a model to classify spam e-mails.
\begin{itemize}
    \item A naive approach is to simply memorize all previous e-mails that have been classified as spam by a human user.
    Then, when a new email arrives, we search for this e-mail in the dataset of all past e-mails and classify accordingly. 
    A problem with this ``\textbf{learning by memorization}'' is that this doesn't generalize well do data outside of the original 
    dataset.

    \item A better approach is to look at the previous set of all emails, an then extract a set of words that correlates strongly
    with classification. Then we apply this rule for all future e-mails. ``\textbf{Inductive reasoning}'' allows us to generalize to data
    outside of the original dataset.

    \item Inductive reasoning may lead us to false conclusions. An example is the \textbf{Pigeon Superstition}. The experiment placed food
    at random intervals initially. Pigeons associated whatever action they did when food arrived and engaged in that activity more.
    Over time, they began to associate the action with the food delivery even though the true pattern was random.

\end{itemize}

\noindent The goal of course, is to develop performant learners. Human learners can filter out this pigeon superstition from common sense,
and we'd desire principles that allow a learner to reject useless conclusions. One powerful tool is the incorporation of \textbf{prior knowledge}
to bias the learning mechanism. The idea is to develop domain specific knowledge, and then translate it into a way to bias the learning
to agree. 


\subsection{Empirical Risk Minimization}
First, we'll lay down some formal notation for learning. In the classical setting:
\begin{itemize}
    \item The learner has access to several datasets for inputs. The \textbf{domain set} \( X \) contains the set of objects to label. Each point may be
    represented using a vector of features. The \textbf{label set} \( y \) contains all possible labels, for now we'll restrict it to a two element set: \( y = \{-1, +1\}   \). 
    Finally, the \textbf{training data} \( S \) consists of a set of \( (x_{m}, y_{m}) \) pairs, the training examples the learner has access to. Namely, 
    \( S = ((x_{1}, y_{1}) \ldots (x_{m}, y_{m})) \).
 
    \item The learner should output a prediction rule \( h \colon X \to y \) that can be used to predict the label of domain points. The function is also known as a 
    \textbf{hypothesis} or a \textbf{predictor}. We also write this hypothesis of some output of our learning algorithm \( A \) upon receiving training sample \( S \).

    \item To generate the data, we sample some points in \( X \) with probability distribution \( D \). Then, we use the true labeling function \( f \colon X \to y \),
    and create labels \( y_{i} = f(x_{i}) \).

    \item The error \( L \) of a classifier is the probability it does not predict the correct label of a random data point over \( D \). Namely,
    
    \[ 
        L_{D, f}(h) \equiv \mathbb{P}_{x \sim D}[h(x) \neq f(x)] \equiv D(\{x: h(x) \neq f(x)\}) 
    \]
    \( L \) is also called the \textbf{generalization error} or the \textbf{risk} or the \textbf{true error}. 
    How we use the distribution D is, if we have some subset \( A \subset X\) then \( D(A) \) gives us how likely it is to observe a point 
    \( x \in A \). 

\end{itemize}

\noindent A good learner should try to minimize the error with respect to some unknown \( D \) and truth \( f \). Because the learner doesn't know \( D \) and \( f \),
we can use a heuristic for the loss called the \textbf{training error}:

\[ 
    L_{S}(h) \equiv \frac{|\{i \in [m] : h(x_{i}) \neq y_{i}\} |}{m} 
\]

\noindent This paradigm, in finding a predictor \( h \) to minimize the training error is \textbf{Empirical Risk Minimization}.


\begin{itemize}
    \item Overfitting is obviously a big problem. A hypothesis that works very well on the training set can generalize very poorly 
    over all data. Additionally, there may exist multiple possible hypothesis that perform equally on the training set, and choosing one is unclear.
    \item A solution is to restrict ERM over a finite set of hypotheses, also called the \textbf{hypothesis class} \( h \in H \). This is a form of 
    inductive bias, but is not always sufficient to guarantee we don't overfit. 
    \item The \textbf{Realizability Assumption} states there exists \( h* \in H \) such that \( L_{(D,f)}(h*) = 0 \). For any random sample \( S \), it's
    possible to find a perfect hypothesis.
    \item The \textbf{i.i.d assumption} assumes the examles in the training set are independently and identically distributed. We denote \( S \sim  D^{m} \)
    where \( D^{m} \) is the probability over m-tuples by applying \( D \) to pick each element of the tuple independently.
\end{itemize}

\noindent Some important metrics to keep in mind are \textbf{confidence} and \textbf{accuracy}. We can denote the probabiltiy of sampling a training set is not 
representative of the underlying \( D \) using the confidence \( (1 - \delta) \). The accuracy parameter can be denoted as \( \epsilon \), where we note a training failure
when \( L_{(D,f)}(h_{S}) \leq \epsilon \)

\begin{itemize}
    \item We'd like to upper bound the probability that some sample will lead to learning failure. Namely, we want to find the bound for
    
    \[ 
        D^{m}(\{S|_{x} \colon L_{(D,f)}(h_{s}) > \epsilon\} ) 
    \]
    \item By the realizability assumption, ERM gives us a hypothesis \( h_{s} \) such that \( L_{S}(h_{s}) = 0\), so the only time learning fails is when 
    our sample is in the set of misleading samples, where there exists hypothesis with 0 training error, but nonzero generalization error. The set of all such samples is also called \( M \), so
    we wish to solve 
    
    \[ 
        D^{m}(\{S|_{x} \colon L_{(D,f)}(h_{s}) > \epsilon\} )  \leq D^{m}(M)
    \]

    \item Since each element is sampled i.i.d, the probability a sampled point is one where a bad hypothesis looks correct is \( 1 - L_{(D,f)}(h) \), which is less than \( 1 - \epsilon \). If we note
    that \( 1 - \epsilon \leq e^{-\epsilon} \) and note that each element is independent, then multiplying the probability over all points in the training set, the probability that a bad hypothesis leads to failure
    over a sample is 
    \[ 
        D^{m}(\{S|_{x} \colon L_{s}(h) = 0\}) \leq (1-\epsilon)^{m} \leq e^{-\epsilon m} 
    \]

    Applying the union bound inequality, we have 
    \[ 
        D^{m}(\{S|_{x} \colon L_{(D,f)}(h_{s}) > \epsilon\} ) \leq |H|e^{\epsilon m}
    \]

    \item For a finite hypothesis class, if we have an integer \( m \) such that \( m\geq\frac{\log{\left( |H|/\delta \right)}}{\epsilon} \), then for any labeling function \( f \) and distribution \( D \)
    where the realizability assumption holds, we have probability of at least \( 1 - \delta \) that we have a bad hypothesis:
    
    \[ 
        Pr_{S\sim D^{m}}[L_{(D,f)}(h_{S}) \leq \epsilon] \geq 1-\delta 
    \]

    In other words, for sufficiently large m, the ERM rule over a finite hypothesis class will be probably (with confidence \( 1-\delta \)) approximately (up to error \( \epsilon \)) correct.

\end{itemize}

\subsection{PAC Learning}




\end{document}